% discussion.tex -- BNR-2026-012

\section{Discussion}
\label{sec:discussion}

\subsection{Hypothesis Verdict}

All four hypothesis components are confirmed with significant margins:

\begin{itemize}
  \item \textbf{H1 (Availability $\geq$ 99.9\%)}: Confirmed. At
    $p = 0.99$, availability is 99.997\% (4.5 nines), exceeding the
    target by 1.5 nines. At $p = 0.999$, availability reaches 7.5 nines.
  \item \textbf{H2 (Attestation $<$ 1\,s)}: Confirmed. P95 at 5\,MB is
    52.7\,ms ($19\times$ margin). P95 at 1\,MB is 10.4\,ms ($96\times$
    margin).
  \item \textbf{H3 (Recovery from 5-of-7)}: Confirmed. 100\% success
    rate at all tested sizes with byte-perfect data reconstruction.
    Recovery with 2 nodes down: 100\% success, $<$1\,ms.
  \item \textbf{H4 (Storage $<$ 2$\times$)}: Confirmed. Total overhead
    is $1.40\times$, matching the theoretical RS minimum of $n/k = 7/5$.
\end{itemize}

\subsection{Privacy Trade-offs}

The hybrid architecture trades information-theoretic data privacy (Shamir
on full data, as in RU-V6) for computational data privacy (AES-256-GCM)
while retaining information-theoretic key secrecy (Shamir on the 32-byte
key). This is a deliberate design choice:

\begin{itemize}
  \item \textbf{Against AES}: if AES-256-GCM is broken (requiring a
    cryptanalytic breakthrough against a NIST-standard cipher), an
    attacker with access to $\geq 5$ RS shards could reconstruct the
    ciphertext and decrypt without the key. However, Shamir key secrecy
    remains intact---the attacker would still need $\geq 5$ key shares.
  \item \textbf{For AES}: AES-256-GCM with 254-bit keys provides
    128-bit security level, considered sufficient for classified
    government communications. The $2.77\times$ storage improvement and
    $36\times$ latency improvement are substantial practical gains.
  \item \textbf{Enterprise context}: for commercial traceability data
    (supply-chain records, production metrics), computational privacy is
    standard practice. Information-theoretic privacy for the data itself
    would be overengineering.
\end{itemize}

\subsection{Key Reduction Modulo BN254}

A critical implementation detail: random 32-byte AES keys can exceed the
254-bit BN254 scalar field prime approximately 75\% of the time. The key
must be reduced modulo $p$ before Shamir sharing. This preserves 254
bits of entropy (the BN254 prime is $\approx 2^{254}$), which exceeds
the 128-bit security level of AES-256-GCM by a factor of $2^{126}$.

Without this reduction, Shamir splitting fails silently or produces
shares that do not reconstruct to the original key. This was identified
and fixed during implementation (see Section~\ref{sec:methodology}).

\subsection{Comparison with Production Systems}

\textbf{vs. EigenDA V2}: EigenDA targets public rollups with 200+
operators and achieves 5\,s average latency. Our 7-node enterprise
committee achieves 4.5\,ms at 500\,KB ($1{,}111\times$ faster), but
this comparison is unfair: EigenDA includes network latency across
hundreds of geographically distributed nodes. The relevant comparison is
architectural: both use RS erasure coding, validating the approach. Our
contribution is the privacy layer (AES+Shamir) absent from EigenDA.

\textbf{vs. Arbitrum AnyTrust}: AnyTrust uses a 6-member committee
with 5-of-6 threshold and BLS aggregation but no erasure coding (each
node stores the full batch). Our $(5,7)$ configuration achieves strictly
higher availability at all reliability levels (Table~\ref{tab:availability-configs})
while storing $1/5$ of the data per node.

\textbf{vs. Polygon CDK DAC}: CDK uses ECDSA multi-sig without erasure
coding or privacy. Our system adds both with negligible latency overhead.

\textbf{vs. RU-V6 (our baseline)}: The hybrid architecture achieves
$36\times$ faster attestation, $2{,}613\times$ faster recovery, and
$2.77\times$ lower storage overhead. The improvement comes from three
sources: (1) Go versus TypeScript execution, (2) SIMD RS versus BigInt
Shamir, and (3) AES-NI hardware encryption.

\subsection{Production Readiness}

The current prototype validates the core protocol but requires four
additions for production deployment:

\begin{enumerate}
  \item \textbf{KZG commitments}: verifiable encoding (INV-DA-P1)
    requires each node to verify its RS chunk against a polynomial
    commitment. KZG verification costs 2 pairing operations
    ($\sim$124{,}000 gas on EVM via BN128 precompiles). Not implemented
    in the prototype.
  \item \textbf{BLS aggregation}: the prototype uses ECDSA P-256
    signatures. Production should use BLS on BN254 for constant-size
    on-chain attestation (48 bytes versus 455 bytes for 7 ECDSA
    signatures).
  \item \textbf{Network latency}: benchmarks are local (no network RTT).
    A 7-node committee with 50--100\,ms RTT per node would add
    100--200\,ms to the attestation pipeline (still within the 1\,s
    target with $\geq 4\times$ margin).
  \item \textbf{On-chain verification contract}: \texttt{BasisDAC.sol}
    must verify attestation certificates (BLS aggregate signature +
    data commitment) before accepting state transitions.
\end{enumerate}

\subsection{Performance Budget with Network Latency}

Table~\ref{tab:budget} estimates the full production latency budget
including network effects.

\begin{table}[t]
\centering
\caption{Estimated Production Latency Budget (500\,KB Batch)}
\label{tab:budget}
\begin{tabular}{@{}lr@{}}
\toprule
\textbf{Operation} & \textbf{Estimated Latency} \\
\midrule
AES-256-GCM encryption              & $\sim$0.05\,ms \\
Shamir key sharing $(5,7)$           & $\sim$0.02\,ms \\
RS encoding $(5{+}2)$               & $\sim$0.06\,ms \\
KZG commitment generation           & $\sim$20\,ms \\
Chunk + key share distribution       & 50--100\,ms \\
KZG chunk verification (per node)    & $\sim$10\,ms \\
ECDSA/BLS signing (per node)         & $\sim$2\,ms \\
Signature collection (parallel)      & 50--100\,ms \\
On-chain verification                & 5--10\,ms \\
\midrule
\textbf{Total estimated}             & \textbf{140--250\,ms} \\
\bottomrule
\multicolumn{2}{l}{\footnotesize Well within 1-second target ($4$--$7\times$ margin).} \\
\end{tabular}
\end{table}

\subsection{Signature Scheme Trade-offs}

Table~\ref{tab:signatures} compares ECDSA and BLS for on-chain
attestation verification.

\begin{table}[t]
\centering
\caption{ECDSA vs BLS for 7-Member Committee}
\label{tab:signatures}
\begin{tabular}{@{}lrr@{}}
\toprule
\textbf{Property} & \textbf{ECDSA (7 sigs)} & \textbf{BLS Aggregated} \\
\midrule
On-chain size        & 455 bytes     & 48 bytes     \\
Verification gas     & 21{,}000     & 124{,}000    \\
Signing latency      & $\sim$1\,ms  & $\sim$2\,ms  \\
Aggregation          & N/A          & $\sim$0.07\,ms \\
Accountability       & Individual   & Bitmap + agg. \\
\bottomrule
\end{tabular}
\end{table}

On the zero-fee Basis Network L1, gas cost is not a constraint. BLS is
preferred for its constant-size signature (48 bytes regardless of
committee size), enabling future committee expansion without increased
on-chain storage.

\subsection{Limitations}

\begin{enumerate}
  \item \textbf{No network latency simulation}: results represent
    computational performance only. Network RTT will dominate the
    attestation pipeline in production.
  \item \textbf{KZG not implemented}: verifiable encoding (INV-DA-P1) is
    specified but not benchmarked. KZG commitment generation may add
    10--50\,ms depending on batch size.
  \item \textbf{ECDSA only}: BLS aggregation is the production target
    but was not benchmarked.
  \item \textbf{Single platform}: benchmarks on Windows 11 AMD64 only.
    Linux deployment may show different AES-NI and SIMD characteristics.
  \item \textbf{Stage 1 only}: stochastic baselines, adversarial
    scenarios (Byzantine nodes, corrupt chunks), and ablation studies
    are planned for subsequent stages.
\end{enumerate}

\subsection{RS vs Shamir: Architectural Comparison}

Table~\ref{tab:rs-vs-shamir} summarizes the fundamental architectural
differences that drive the performance gains.

\begin{table}[t]
\centering
\caption{Reed-Solomon vs Shamir SSS for $(5,7)$ DAC}
\label{tab:rs-vs-shamir}
\begin{tabular}{@{}lcc@{}}
\toprule
\textbf{Property} & \textbf{Shamir $(5,7)$} & \textbf{Hybrid (AES+RS)} \\
\midrule
Privacy model     & Info-theoretic   & Computational + IT key \\
Per-node storage  & $1\times$ original & $1/5$ original \\
Total storage     & $7\times$        & $1.4\times$ \\
Encoding time     & $O(kn)$ per elem & $O(n \log n)$ total \\
Decoding time     & $O(k^2)$ per elem & $O(k \log^2 k)$ total \\
Verifiable        & No               & With KZG: yes \\
Key management    & None             & AES key distribution \\
\bottomrule
\end{tabular}
\end{table}
